\chapterhead{4}{Machines Trade Distance}{Less push, but over farther}\label{ch4}

A \textbf{lever} is a stiff bar that turns on a point called the pivot. A \textbf{pulley} is a wheel with a rope over it. A \textbf{ramp} is a slope. These simple machines do not make energy. They let you trade: you push with less force, but you push over a longer distance.

For the lever the trade is easy to see. Multiply the number of coins by their distance from the pivot on each side. When the two products match, the lever balances.

\begin{center}
\sffamily\small\renewcommand{\arraystretch}{1.5}
\begin{tabular}{@{}p{3.4cm}p{3.7cm}p{3.7cm}@{}}
\toprule
\textbf{Model} & \textbf{What it keeps} & \textbf{What it drops}\\
\midrule
Coins $\times$ distance & the balance rule & the weight of the ruler\\
Pull $\times$ distance & the trade a pulley makes & friction in the wheel\\
A machine as a trade & less force for more distance & everything inside it\\
\bottomrule
\end{tabular}
\end{center}

\begin{tryit}{Balance a ruler}
\textbf{You need:} a 30~cm ruler, a pencil, coins of one kind and sticky tape.
\begin{enumerate}[leftmargin=1.6em]
\item Rest the ruler on the pencil at its 15~cm mark (the pivot).
\item Stack 6 coins 2~cm from the pivot on one side. That is $6\times2=12$.
\item Find where 3 coins balance, then 2 coins, then 1 coin. Write each distance and the product.
\end{enumerate}
\par\medskip
\renewcommand{\arraystretch}{1.6}
\begin{tabular}{|l|c|c|c|}
\hline
\textbf{Coins on far side} & \textbf{3} & \textbf{2} & \textbf{1}\\ \hline
\textbf{Balance distance (cm)} & & & \\ \hline
\textbf{Coins $\times$ distance} & & & \\ \hline
\end{tabular}
\end{tryit}

\begin{worked}{Practice numbers, not from a real test}
Load: 6 coins at 2~cm, so 12. Three coins balanced at 4~cm, giving $3\times4=12$. Two coins balanced at 6~cm, giving 12. One coin balanced at 12~cm, giving 12.\\
Fewer coins needed more distance. The product stayed the same. Real balances wobble by a centimetre or so, so real products may differ slightly.
\end{worked}

\begin{tryit}{A pulley}
Hang a bag of coins from a string over a pencil or a spool, and lift it 10~cm by pulling straight down. Use a rubber band in the line and write how far it stretches. Then loop the string under a second spool attached to the bag, fix one end above and pull up on the other end. Write the stretch and how far you had to pull the string.\par\medskip
\renewcommand{\arraystretch}{1.6}
\begin{tabular}{|l|c|c|c|}
\hline
 & \textbf{Stretch (cm)} & \textbf{String pulled (cm)} & \textbf{Stretch $\times$ string}\\ \hline
\textbf{Straight} & & 10 & \\ \hline
\textbf{With pulley} & & & \\ \hline
\end{tabular}
\end{tryit}

\begin{worked}{Practice numbers}
Straight: stretch 6, string 10, product 60. Pulley: stretch 3.2, string 20, product 64.\\
The pulley halved the pull but needed twice the string. The product went from 60 to 64. The extra 4 is a gap, and friction in the wheel is a likely place: it turned some motion into warmth.
\end{worked}

\begin{multix}[hmm]\multixbot[hmm]\tcblower
A lever needs only 2 coins instead of 6. That is less energy overall, so a lever saves energy. I will lever everything. Is that right?
\end{multix}

\trust{We multiplied each side of the lever and each pull by its distance, and we compared the products.}
{On the lever the products matched. On the pulley the products were close, and the extra was named as friction.}
{That a machine saves energy. It saves effort for a moment and costs distance.}

\begin{keepbox}
We turned a machine into a product. \textbf{We kept} \wline[3.4cm] \textbf{and we set aside} \wline[3.4cm].
\end{keepbox}

\clearpage

\begin{puzzle}{Balance}
6 coins sit at 2~cm. Where do 4 coins balance? \wline[1.6cm]~cm
\end{puzzle}

\begin{puzzle}{Another}
8 coins sit at 3~cm. Where do 2 coins balance? \wline[1.6cm]~cm
\end{puzzle}

\begin{puzzle}{Pulley}
Lifting straight takes a pull of 5 over 10~cm, a product of 50. A pulley halves the pull. How far is the string pulled? \wline[1.6cm]~cm
\end{puzzle}

\begin{puzzle}{Fit the job}
Write L for lever, R for ramp, P for pulley or W for wheel.\\
(a) raise a flag \wline[1cm] \quad (b) roll a barrel into a truck \wline[1cm] \quad (c) lift a rock a little way \wline[1cm] \quad (d) move a cart \wline[1cm]
\end{puzzle}

\begin{puzzle}{The trade}
A lever has 8 coins at 5~cm. A helper pushes 20~cm out. How many coins' worth of push is needed? \wline[1.6cm] How many times farther does the helper push? \wline[1.6cm]
\end{puzzle}

\clearpage

\begin{puzzle}{Spot the cheat}
A diagram says a pulley lifts a load 10~cm with half the pull and only 10~cm of string. What is wrong? \wline[6cm]
\end{puzzle}

\begin{puzzle}{Gap}
Straight pull product 60. Pulley product 64. Gap \wline[1.6cm]. Name a likely cause. \wline[4cm]
\end{puzzle}

\begin{puzzle}{Draw the circles}
A pulley takes in 64 units. 60 lift the bag. Fill in the circle for the rest.
\begin{center}
\begin{tikzpicture}
  \node[bigcir,minimum size=2.4cm] (a) at (0,0) {}; \node[font=\sffamily\scriptsize] at (0,0) {pull 64};
  \node[cir] (b) at (4,0.8) {60}; \node[cir,dashed] (c) at (4,-0.8) {};
  \draw[popA] (a) -- (b); \draw[popA] (a) -- (c);
\end{tikzpicture}
\end{center}
\end{puzzle}

\begin{puzzle}{Wrap}
A student wrote ``A lever makes work for free.'' Wrap it with a better claim. \wline[6cm]
\end{puzzle}

\begin{puzzle}{Design a lifter}
Plan a machine that lifts a heavy book with a light push. Say what you would measure. \wline[6cm]
\end{puzzle}
\clearpage

\thispagestyle{fancy}
\subhead{Lever record}
Try three loads on the same ruler. Keep the pivot at the middle.\par\medskip
\renewcommand{\arraystretch}{3.0}
\begin{tabularx}{\linewidth}{|>{\sffamily\small\bfseries}p{2.6cm}|X|X|X|X|}
\hline
Load & Distance (cm) & Balancing coins & Distance (cm) & Product match? \\ \hline
 & & & & \\ \hline
 & & & & \\ \hline
 & & & & \\ \hline
\end{tabularx}

\vspace{8pt}
Which balance was easiest to find? \wline[7cm]\\[10pt]
Did any product differ, and by how much? \wline[5cm]\\[10pt]
Sketch your lever with its coins and distances.\par\medskip
\begin{tikzpicture}\draw[ink!70,line width=1pt,rounded corners=3pt] (0,0) rectangle (11.4,5.2);\end{tikzpicture}
\clearpage
